\section{Parallel Validation}
\label{sec:parallel}

A partitioned second-order solver must give the same answer regardless of how the
mesh is cut. The rank-count study covers one NACA case and one cylinder case at
$np=1,2,4,8$, so that both meshes and both wall types are exercised, and
\cref{tab:scaling} reports the resulting forces, edge cuts, and timings.

\begin{table}[htbp]
  \centering
  \caption{Rank-count study. Each row is an independent run of the same case at a
  different rank count, with the final $C_D$ taken from \texttt{forces.csv} and the
  edge cut from \texttt{metadata.json}.}
  \label{tab:scaling}
  \small
  \begin{tabular}{lrrrrr}
\toprule
Case & Ranks & Edge cut & $C_D$ at step 1500 & rel.\ dev.\ from $np{=}1$ & Wall (s)$^\ddagger$ \\
\midrule
cylinder\_m010\_laminar\_re20 & \num{1} & \num{0} & 2.13640805 & \num{0.0e+00} & \num{24.32} \\
cylinder\_m010\_laminar\_re20 & \num{2} & \num{97} & 2.13670466 & \num{1.4e-04} & \num{12.27} \\
cylinder\_m010\_laminar\_re20 & \num{4} & \num{282} & 2.13691268 & \num{2.4e-04} & \num{6.258} \\
cylinder\_m010\_laminar\_re20 & \num{8} & \num{471} & 2.13666609 & \num{1.2e-04} & \num{8.883} \\
naca0012\_m015\_laminar\_re5000 & \num{1} & \num{0} & 0.11886373 & \num{0.0e+00} & \num{46.41} \\
naca0012\_m015\_laminar\_re5000 & \num{2} & \num{120} & 0.11886855 & \num{4.1e-05} & \num{21.69} \\
naca0012\_m015\_laminar\_re5000 & \num{4} & \num{241} & 0.11887621 & \num{1.1e-04} & \num{10.48} \\
naca0012\_m015\_laminar\_re5000 & \num{8} & \num{439} & 0.11886993 & \num{5.2e-05} & \num{13.76} \\
\bottomrule
\end{tabular}

\end{table}

\subsection{Consistency across rank counts}

The measured consistency is that $C_D$ agrees to \textbf{five significant
figures} across $np=1,2,4,8$, with the METIS edge cut rising through
$0, 97, 282, 471$ as the partition is refined. Agreement at this level is the
expected and correct result, and it is worth being precise about why it is not
exact.

Three mechanisms could make results rank-dependent, and two of them are
eliminated by construction:

\begin{itemize}
  \item \textbf{Reconstruction across a cut.} A partition face must reconstruct
  identically from both sides. This holds because the halo exchange carries
  gradients and limiter factors as well as the state (\cref{sec:mpi}), so
  \cref{eq:recon} sees identical inputs on both ranks. Without the gradient and
  limiter exchange, the scheme would silently degrade to first order at every cut
  and the drag would change with rank count at the percent level.
  \item \textbf{Force double-counting.} A wall face on a partition boundary must
  be integrated exactly once. This is enforced by integrating only faces whose
  adjacent cell is owned (\cref{sec:forces}).
  \item \textbf{Floating-point summation order.} This one is irreducible. The
  global residual and force reductions sum contributions in an order that depends
  on the partition, and floating-point addition is not associative, so results
  differ in the last few digits. Those differences then feed back through the
  nonlinear iteration, which is why agreement is to five figures rather than to
  machine precision.
\end{itemize}

The remaining rank dependence is in the \emph{path} to the solution rather than
the solution itself. As noted in \cref{sec:lusgs}, LU-SGS is exactly Gauss--Seidel
within a rank and block-Jacobi across ranks, because off-rank increments are
frozen between sweeps. Increasing the rank count therefore weakens the implicit
coupling slightly and can change the iteration count needed to reach a given
residual level, without changing the converged answer. This is a well-known and
accepted property of LU-SGS in parallel, and it is the reason the inner-iteration
counts in \cref{tab:scaling} need not match across rank counts even though the
forces do.

\subsection{Load balance and communication}

\Cref{tab:partition} gives the per-rank decomposition for a representative $np=8$
run, with a load-balance ratio of $\cnsPartitionBalance$. METIS balances the cell
count well, but two caveats apply to interpreting that number as
time-to-solution balance:

\begin{itemize}
  \item Boundary faces are not distributed in proportion to cells. A rank holding
  part of the airfoil surface does more work per cell than a rank holding only
  farfield cells, because wall faces carry the extra boundary-state and
  viscous-correction work of \cref{eq:wallgrad}. The
  \texttt{num\_boundary\_faces} column of \cref{tab:partition} shows this spread
  directly.
  \item Neighbour counts vary between ranks, so the halo-exchange cost is not
  uniform. The interior ranks of the partition have the most neighbours and thus
  the largest message counts per exchange.
\end{itemize}

Communication volume per exchange is small: each neighbour pair moves the cells
listed in the send/receive columns of \cref{tab:partition}, times 4 components for
the state or limiter and 8 for the gradients. Because the exchanged surface scales
as the partition perimeter while the work scales as the partition area, the
communication fraction grows with rank count on a fixed mesh; on meshes of this
size ($\num{20816}$ and $\num{10185}$ cells) the $np=8$ runs are already in the
regime where the per-rank subdomain is small enough that this overhead is visible
in the timings of \cref{tab:scaling}.

\subsection{Choice of rank count for production}
\label{sec:rankchoice}

These meshes are small enough that more ranks are not better, and the production
rank count was chosen from measurement rather than by assumption. Timing a fixed
20 physical steps of the transient cylinder case gave $11.3$~s at $np=8$,
$12.6$~s at $np=16$ and $26.4$~s at $np=32$: past eight ranks the solver gets
monotonically \emph{slower}. With only $\cylCells$ cells, an $np=32$ partition
leaves roughly 300 cells per rank, at which point the halo exchange and the
serialization inherent in the LU-SGS sweep dominate the useful work per sweep.
All production cases therefore run at $np=8$, which is also the rank count the
benchmark requires to be demonstrated.

One measurement artefact is worth recording so the timings in this report are
interpreted correctly. An earlier attempt to run several cases concurrently at
high rank counts oversubscribed the machine --- 69 MPI processes on 64 cores, each
receiving roughly $5\%$ of a core --- and the resulting wall times are meaningless
as solver performance. All timings reported here are from runs made after that
configuration was corrected, on an otherwise unloaded machine. Wall-clock
comparisons between cases are only meaningful under that condition, which is why
it is stated rather than assumed.

\section{Visualization and Plot Style Requirements}
\label{sec:visualization}

All figures are generated by a single reproducible pipeline
(\texttt{tools/make\_figures.py}) reading only submitted output files, and every
figure is recorded in \texttt{report/figure\_manifest.csv} together with the case,
figure type, plotted variable, source file, and caption. The manifest is what makes
each figure traceable to its data.

The plotting conventions are uniform: line plots use readable font sizes,
consistent line widths and a colour-blind-safe palette, labelled and
nondimensionalised axes, legends, and a light grid; residual histories use a
logarithmic ordinate. Field plots are filled contour renderings on the actual
unstructured triangulation --- quadrilaterals are split for rendering only, never
for the solution --- rather than scatter point clouds, and each carries a colourbar
labelled with its variable name. Near-body views are used for the airfoil and
cylinder so the boundary layer, shocks and wake are visible at a useful scale, with
separate whole-domain views where the farfield behaviour matters.

Two conventions are worth stating explicitly because they are contract
requirements. First, file names match plotted variables exactly: a file named
\texttt{..\_mach.png} plots Mach number and \texttt{..\_pressure.png} plots
pressure, and both are produced for every case. Second, vorticity plots use the
recommended symmetric clipped range $[-5,5]$; the near-wall vorticity of a no-slip
cylinder is more than an order of magnitude larger than the wake vorticity, so
automatic scaling renders the vortex street invisible.
